% Generalized Phase Projection Operator P_phi
% Version: 1.0
% Date: 2026-03-03
% Status: Canonical - GEOMETRY
% Author: Ing. David Jaroš
% © 2026 Ing. David Jaroš — CC BY-NC-ND 4.0

\section{Generalized Phase Projection $P_\varphi$}
\label{sec:canonical:phase_projection}

\subsection{Definition}

Let $\mathcal{G}_{\mu\nu} \in \mathbb{C} \otimes \mathbb{H}$ be the fundamental biquaternionic metric
(Axiom C, \texttt{core/AXIOMS.md}).  For any constant phase angle $\varphi \in \mathbb{R}$ we define
the \textbf{generalized real projection}

\begin{equation}
\label{eq:phase_projection_def}
\boxed{P_\varphi[\mathcal{G}_{\mu\nu}] \;:=\; \mathrm{Re}\!\left(e^{-i\varphi}\,\mathcal{G}_{\mu\nu}\right).}
\end{equation}

Because $\mathcal{G}_{\mu\nu}$ takes values in the complex–quaternionic algebra $\mathbb{C}\otimes\mathbb{H}$,
the factor $e^{-i\varphi}$ acts by the standard complex scalar multiplication, and $\mathrm{Re}(\cdot)$
extracts the real part of each complex coefficient in the biquaternionic expansion.

The result $P_\varphi[\mathcal{G}_{\mu\nu}]$ is a real-valued, symmetric $4\times 4$ tensor and therefore
a valid Lorentzian metric tensor on the same underlying manifold.

\subsection{Recovery of the Standard Observer Metric ($\varphi = 0$)}

Setting $\varphi = 0$:
\begin{equation}
\label{eq:phase_proj_phi0}
P_0[\mathcal{G}_{\mu\nu}] = \mathrm{Re}\!\left(e^{0}\,\mathcal{G}_{\mu\nu}\right) = \mathrm{Re}(\mathcal{G}_{\mu\nu}) = g_{\mu\nu}.
\end{equation}

This is precisely the unique emergent metric of Axiom C.  Hence $\varphi = 0$ is the
\textbf{standard observer (gauge) choice}, not a theoretical requirement.  The general
framework uses $P_\varphi$ for any constant $\varphi$.

\begin{tcolorbox}[colback=blue!5!white,colframe=blue!75!black,title=Note on Axiom C]
The formula $g_{\mu\nu} = \mathrm{Re}(\mathcal{G}_{\mu\nu})$ stated in Axiom C is the \emph{special case}
$\varphi = 0$ of the general operator $P_\varphi$.  Axiom C is unchanged; the present section
formalizes the full one-parameter family of which that definition is a member.
\end{tcolorbox}

\subsection{Linearity}

$P_\varphi$ is a $\mathbb{R}$-linear operator on $\mathbb{C}\otimes\mathbb{H}$-valued tensors:
\begin{equation}
P_\varphi[\alpha\,\mathcal{A}_{\mu\nu} + \beta\,\mathcal{B}_{\mu\nu}]
= \alpha\,P_\varphi[\mathcal{A}_{\mu\nu}] + \beta\,P_\varphi[\mathcal{B}_{\mu\nu}],
\qquad \alpha,\beta\in\mathbb{R}.
\end{equation}

\subsection{Commutativity with Partial Derivatives ($\varphi = \mathrm{const}$)}

When $\varphi$ is constant across spacetime,

\begin{equation}
\label{eq:proj_commutes_partial}
\partial_\mu P_\varphi[\mathcal{G}_{\nu\rho}]
= P_\varphi[\partial_\mu \mathcal{G}_{\nu\rho}],
\end{equation}

because $e^{-i\varphi}$ is a constant scalar and $\partial_\mu$ commutes with $\mathrm{Re}(\cdot)$ for
smooth fields.  This commutativity extends to the gauge-covariant derivative $D_\mu$ when $\varphi$ is
covariantly constant ($\nabla_\mu\varphi = 0$), a property used in
Section~\ref{sec:canonical:gr_as_limit}.

\subsection{Preservation of Tensor Dimensionality}

The map $P_\varphi$ acts coefficient-by-coefficient on the $4\times 4$ index structure of
$\mathcal{G}_{\mu\nu}$; it does not mix spacetime indices.  Therefore
$P_\varphi[\mathcal{G}_{\mu\nu}]$ is a $(0,2)$-tensor of the same rank and dimensionality as the
classical metric $g_{\mu\nu}$ for every value of $\varphi$.

\subsection{Phase Frame Bundle}

The collection of all phase observers is parametrized by $\varphi \in \mathrm{U}(1) \subset \mathbb{C}$.
We define the \textbf{phase frame bundle}

\begin{equation}
\mathcal{F} \;=\; \mathcal{M} \times_{\varphi} \mathrm{U}(1),
\end{equation}

where $\mathcal{M}$ is the four-dimensional spacetime manifold.  Each fiber $\{(\varphi)\}_{\varphi\in\mathrm{U}(1)}$
labels a \emph{phase frame}, and the section $\varphi = \mathrm{const}$ defines a \emph{constant phase
observer}.  The standard observer uses the section $\varphi = 0$.

The projection $P_\varphi$ is the canonical projection map from the total space of $\mathcal{F}$ to the
space of real Lorentzian metrics.

\subsection{KK Mode Expansion and the Ultrametric Structure of Phase States}
\label{subsec:kk_ultrametric}

\subsubsection{Kaluza--Klein Mode Expansion}

Decompose the $\Theta$-field in the imaginary-time direction $\psi$ (Axiom B):

\begin{equation}
\label{eq:kk_expansion}
\Theta(x, \psi) = \sum_{n=-\infty}^{\infty} a_n(x)\,e^{in\psi},
\end{equation}

where $a_n(x) \in \mathbb{C}\otimes\mathbb{H}$ are biquaternionic mode amplitudes and $n \in \mathbb{Z}$
is the KK mode number.

\subsubsection{Prime-Indexed Stable Modes}

The drift-diffusion equation governing phase dynamics (Appendix~H,
\texttt{Appendix\_H\_Theta\_Phase\_Emergence.tex}) admits attractor solutions at mode numbers
corresponding to \textbf{prime integers}:

\begin{equation}
n \in \mathcal{P} = \{2, 3, 5, 7, 11, \ldots\},
\end{equation}

because prime-mode perturbations are maximally isolated from resonant coupling — no proper
divisors other than $1$ and $n$ itself contribute to mixing terms in the nonlinear attractor equation.
All other modes are driven to zero amplitude in the long-time limit.

\subsubsection{P-adic Metric on Phase States}

For each prime $p$, the \textbf{$p$-adic absolute value} $|\cdot|_p$ on $\mathbb{Q}$ satisfies the
ultrametric inequality $|x + y|_p \leq \max(|x|_p, |y|_p)$.  We define the
\textbf{$p$-adic distance} on the phase-state space by

\begin{equation}
\label{eq:padic_distance}
d_p(\varphi_1, \varphi_2) \;:=\; |\varphi_1 - \varphi_2|_p.
\end{equation}

This distance satisfies the \textbf{ultrametric inequality}:
\begin{equation}
d_p(\varphi_1, \varphi_3) \leq \max\!\bigl(d_p(\varphi_1, \varphi_2),\, d_p(\varphi_2, \varphi_3)\bigr).
\end{equation}

For a rigorous treatment of $p$-adic analysis applied to physics see the existing p-adic
material in \texttt{consolidation\_project/appendix\_ALPHA\_padic\_derivation.tex}.

\subsubsection{Orthogonality of Prime-Indexed Modes}

Under the $p$-adic metric, prime-indexed modes $e^{in\psi}$ and $e^{im\psi}$ with distinct primes
$n \neq m$ satisfy
\begin{equation}
d_n(n, m) = |n - m|_n = 1,\quad d_m(n, m) = |n - m|_m = 1.
\end{equation}
Since no prime divides another, the two modes live in $p$-adically non-overlapping
balls, making them \textbf{orthogonal in the $p$-adic sense}: they cannot be continuously
deformed into each other without passing through all intermediate integers.

\subsubsection{Ultrametric Space of Stable Phase States}

\begin{tcolorbox}[colback=green!5!white,colframe=green!75!black,title=Theorem: Ultrametric Phase Space]
\textbf{The set of stable phase states of UBT forms an ultrametric space under the $p$-adic metric.}

More precisely: the stable $\varphi$ values are indexed by prime KK mode numbers
$\{p \in \mathcal{P}\}$, and the $p$-adic distance $d_p(\varphi_i, \varphi_j)$ endows this
set with an ultrametric topology distinct from the standard Archimedean metric on $\mathbb{R}$.
\end{tcolorbox}

\subsubsection{Connection to the Fine Structure Constant}

The prime $p = 137$ is a candidate stable mode because $\alpha \approx 1/137$ and $137$ is prime.
In the KK expansion, the mode $n = 137$ sits at the boundary of the non-perturbative regime
($n \sim \alpha^{-1}$), making it a natural attractor of the phase dynamics.  A detailed derivation
is given in \texttt{consolidation\_project/appendix\_ALPHA\_padic\_derivation.tex}.

% End of generalized phase projection
